\documentclass{article}
\usepackage{microtype,graphicx,booktabs,amsmath,amssymb,url}
\usepackage{hyperref}
\usepackage[accepted]{icml2025}
\usepackage{xcolor}

\icmltitlerunning{Repairing the Prior of a Tabular Foundation Model}

\input{numbers}

\begin{document}

\twocolumn[
\icmltitle{Repairing the Prior: An LLM Agent that Revises the Synthetic\\
Pre-training Distribution of a Tabular Foundation Model}
\icmlsetsymbol{equal}{*}
\begin{icmlauthorlist}
\icmlauthor{Anonymous}{anon}
\end{icmlauthorlist}
\icmlaffiliation{anon}{Anonymous Institution}
\icmlcorrespondingauthor{Anonymous}{anon@example.com}
\icmlkeywords{tabular foundation models, in-context learning, LLM agents, synthetic priors}
\vskip 0.3in
]

\printAffiliationsAndNotice{}

\begin{abstract}
Tabular foundation models (TFMs) such as TabPFN and TabICL are pre-trained on millions of
synthetic datasets drawn from a hand-designed prior over structural causal models. That
prior is frozen at publication time and knows nothing about the tasks the model will
actually be asked to solve; recent work argues prior design, not architecture, is now the
binding constraint on TFM quality. Re-pre-training a TFM to fix its prior costs a GPU
cluster, so in practice nobody does.

We show the prior can instead be \emph{repaired}, cheaply and after the fact. We treat the
prior's generating configuration as a search space and give it to an LLM agent, which reads
a diagnosis of where the released checkpoint fails on a small set of development tasks and
proposes a revised prior. Each candidate prior is materialised, the released checkpoint is
LoRA-adapted to it ($2.1\%$ of parameters, \LoRAMinutes{} on one A100), and the result is
scored. The winning adapter is then evaluated on \NTest{} held-out tasks the agent never
saw.

On real OpenML tables the agent-revised prior changes held-out AUC by \APMean{}
(bootstrap $95\%$ CI $[\APCIlo{}, \APCIhi{}]$) relative to the released checkpoint, and by
\ARMean{} (CI $[\ARCIlo{}, \ARCIhi{}]$) relative to a random search over the same knobs given
an identical budget --- the comparison that separates the LLM's contribution from the search
itself. Inspecting the winning configurations shows what the agent found: TabICL's prior
over-represents large, many-class, purely numeric datasets, and the agent shifts it toward
the small, binary, partly categorical regime that real tables actually occupy.
\end{abstract}

\section{Introduction}

A tabular foundation model does not learn from tables. It learns from a \emph{prior}: a
program that samples synthetic datasets, typically by drawing a random structural causal
model (SCM), pushing noise through it, and reading off features and labels. TabPFN
\cite{hollmann2023tabpfn,hollmann2025tabpfnv2} and TabICL \cite{qu2025tabicl} are both
trained this way, and the resulting model performs Bayesian in-context inference over
whatever prior it was shown. The prior \emph{is} the inductive bias.

This has an awkward consequence. The prior is chosen once, by the model's authors, before
any downstream task exists. If it over-represents regimes that real tables never occupy ---
too many classes, too many rows, too few categorical columns --- the model wastes capacity
on a world it will never see, and the mismatch is invisible at pre-training time because no
real task is in the loop. The field has begun to notice: \citet{zhang2025mitra} argue that
``more attention should be paid to the design of the data prior'' and introduce the first
systematic prior-quality criteria, while O'PRIOR \cite{bouadi2026oprior} shows that holding
architecture and compute fixed and varying only the prior moves downstream accuracy
substantially.

The obvious fix --- design a better prior and pre-train again --- is the one nobody can
afford. Pre-training a TFM consumes millions of synthetic datasets on a GPU cluster. So the
prior stays broken.

\paragraph{This paper.} We show the prior can be repaired \emph{post hoc}, without
re-pre-training, by an agent that is allowed to see downstream failure. Our contributions:

\begin{itemize}\itemsep1pt
\item We first show the mismatch is \emph{real and measurable}, with no model in the loop:
sampling from TabICL's released prior and measuring it on the same axes as real OpenML tables
shows it generates tables three times wider, far more class-imbalanced, with almost no feature
redundancy and no categorical columns (Table~\ref{tab:audit}).
\item We recast prior design as a \emph{closed-loop search problem}. The prior's generating
configuration --- SCM depth and width, noise scale, class count, sequence length,
categorical fraction, MLP/tree mixture --- is exposed as an editable object. An LLM agent
proposes revisions conditioned on a diagnosis of where the released checkpoint currently
fails, and never sees a real row or even a task's name.
\item We make each candidate prior cheap to evaluate. Rather than pre-training from
scratch, we LoRA-adapt the released checkpoint to the candidate prior, touching $2.1\%$ of
parameters. One candidate costs \LoRAMinutes{} on a single A100, turning prior design from a
cluster job into a search loop.
\item We isolate the LLM's contribution with a random-search arm that draws priors uniformly
from the same configuration space under the same budget. This is the baseline that decides
whether world knowledge helps or whether any search would have done.
\item We report the outcome on \NTest{} held-out real tasks as paired differences with
confidence intervals, and we show \emph{what} the agent changed --- the revisions are
interpretable and consistent across seeds, which is itself evidence that the mismatch it
found is real rather than an artefact of search.
\end{itemize}

\section{Method}

\paragraph{Setup.} Let $M_{\theta_0}$ be a released TFM checkpoint, pre-trained on a prior
$P_0$ with configuration $c_0$. A prior configuration $c$ induces a distribution over
synthetic datasets; sampling from it is cheap, pre-training on it is not. We are given a
small set of \emph{development} tasks $\mathcal{D}$ (real tables, which the agent may
inspect) and a disjoint set of \emph{test} tasks $\mathcal{T}$, which nothing in the loop
ever sees.

\paragraph{The knob space.} We expose \Nknobs{} configuration knobs of TabICL's SCM prior,
each with hard bounds (Table~\ref{tab:knobs} lists those the agent moved). They control the
causal graph (number of causes, SCM depth, hidden width, initialisation and noise scale),
the sampling regime (dataset size, feature count, class count, train/test split, whether
small datasets are replayed), and the mixture between MLP-SCM and tree-SCM generators.
Proposals outside the bounds are clipped, so a bad proposal costs a training run but cannot
produce an invalid prior.

\paragraph{Cheap evaluation via LoRA.} Given a candidate $c$, we instantiate its prior,
sample synthetic tasks, and adapt $\theta_0$ by LoRA \cite{hu2022lora} on the feed-forward
projections of the dataset-wise in-context transformer --- the block that performs the
actual in-context inference. The column and row embedders keep their pre-trained behaviour.
This trains \LoRAParams{} of \TotalParams{} parameters ($2.1\%$) with the standard
in-context objective: sample a synthetic dataset, split it, condition on the training part,
cross-entropy on the held-out part. The adapted model is then scored on $\mathcal{D}$
exactly as a user would score it --- fit the context, one forward pass, AUC.

\paragraph{The agent.} At round $r$ the agent receives (i) a card describing each
development task --- feature count, class count, context size, class balance --- (ii) the
per-task AUC of the current adapter, worst first, and (iii) the search history: every
configuration tried and its development score. It returns a full revised configuration. The
loop keeps the configuration with the best development score; that adapter, and only that
one, is finally evaluated on $\mathcal{T}$.

Two properties of this design matter for what we can claim. First, the development tasks are
presented \emph{anonymously}: the agent sees summary statistics, never raw rows, and never
the task's name. This is deliberate. LLMs have memorised the contents of well-known OpenML
tables \cite{bordt2024elephants}, so naming them would let recalled specifics --- rather than
reasoning about what real tabular problems look like --- drive the revision. Second, the
agent never sees the test tasks at all. What transfers, therefore, is a belief about the
shape of real tabular problems, not knowledge of any particular one.

\paragraph{Baselines with identical budget.} The arm that decides the paper is
\textbf{random search}: the same number of candidate priors, drawn uniformly from the same
bounded knob space, the same LoRA compute per candidate, the same
best-on-development selection rule. The \emph{only} difference is who proposes the next
configuration. We also report \textbf{base}, which LoRA-adapts to TabICL's own unmodified
prior --- this controls for the possibility that any further training on any prior helps.

\section{Experiments}
\label{sec:exp}

\paragraph{Data.} \NDev{} OpenML tasks form the development set and \NTest{} disjoint tasks
form the test set, spanning both semantically named tables (\texttt{cmc}, \texttt{churn},
\texttt{spambase}) and anonymised ones (\texttt{ozone}, \texttt{phoneme},
\texttt{bank-anon}). Each task is cut into a $200$-row context, a validation split, and a
test split; we report AUC (macro one-vs-rest for multiclass) of the frozen model conditioned
on the context. We deliberately exclude \texttt{adult}, which \citet{bordt2024elephants} show
is memorised verbatim by GPT-class models.

\paragraph{Protocol.} Backbone: released \texttt{tabicl-classifier-v2}. \NRounds{} prior
candidates per search arm, \Steps{} LoRA steps each, \NSeeds{} seeds. All arms share one
A100.

\paragraph{The mismatch is measurable before any agent runs.}
Our premise is that the released prior does not look like real tabular data. This is
checkable directly, without any model: sample datasets from the prior, take real OpenML
tables, and measure both on the same axes (Table~\ref{tab:audit}). The released prior
generates tables with a median of \PriorFeat{} features against \RealFeat{} in real data,
class distributions far more skewed than real ones (minority fraction \PriorMinority{} vs
\RealMinority{}), features that are nearly uncorrelated where real columns are redundant, and
essentially no low-cardinality (categorical-like) columns. The prior is not a bad prior --- it
is a prior for a different population of tables than the one it meets.

\input{audit}

\begin{table*}[t]
\centering
\small
\caption{Held-out AUC on \NTest{} test tasks the agent never saw. All rows use the same
frozen TabICLv2 backbone; the only difference is the prior its LoRA adapter was trained on.
Random search and the agent are given an identical budget (\NRounds{} candidate priors,
\Steps{} LoRA steps each), so the gap between them is attributable to the proposer.
Mean $\pm$ s.d.\ over \NSeeds{} seeds.}
\label{tab:main}
\begin{tabular}{lccccccc}
\toprule
& \multicolumn{6}{c}{test task} & \\
\cmidrule(lr){2-7}
prior the adapter was trained on & cmc & churn & spambase & ozone & phoneme & bank-anon & mean \\
\midrule
\maintablebody
\bottomrule
\end{tabular}
\end{table*}

\paragraph{Main result.} Table~\ref{tab:main} reports held-out AUC. Because arms share tasks
and seeds, we analyse \emph{paired} differences --- one per (task, seed) pair, $n=\ARN{}$ ---
rather than three seed means, and report a bootstrap $95\%$ CI and a Wilcoxon signed-rank
$p$.

Adapting to TabICL's own unmodified prior changes held-out AUC by \BPMean{}
(CI $[\BPCIlo{}, \BPCIhi{}]$): simply doing more LoRA training on the prior the model was
already pre-trained on does not buy anything. Searching over priors does. Random search
gains \RPMean{} (CI $[\RPCIlo{}, \RPCIhi{}]$) over the released checkpoint, and the
agent-revised prior gains \APMean{} (CI $[\APCIlo{}, \APCIhi{}]$, $p=\APP{}$).

The comparison the paper turns on is agent \emph{versus} random search, which holds the
budget, the knob space, the LoRA compute and the selection rule fixed and varies only who
proposes the next prior: \ARMean{} (CI $[\ARCIlo{}, \ARCIhi{}]$, $p=\ARP{}$), winning on
\ARWins{} of \ARN{} (task, seed) pairs.

\paragraph{Search efficiency.} Figure~\ref{fig:search} plots best-so-far development AUC
against the number of prior candidates evaluated. The agent's advantage is concentrated in
the first candidates: it does not need to stumble onto a good region of the knob space,
because it can reason about the mismatch directly.

\begin{figure}[t]
\centering
\includegraphics[width=\columnwidth]{fig_search.pdf}
\caption{Best development AUC as a function of prior candidates evaluated (mean $\pm$ s.d.\
over seeds). Both search arms pay the same LoRA compute per candidate.}
\label{fig:search}
\end{figure}

\paragraph{What did the agent change?} The revisions are interpretable and consistent across
seeds (Table~\ref{tab:knobs}). The agent narrows \texttt{max\_classes} from TabICL's $10$
toward $2$, having observed that every development task is binary or near-binary; shrinks
\texttt{max\_seq\_len} and enables \texttt{replay\_small}, matching the small-table regime
TFMs are actually deployed in; and raises \texttt{cat\_prob}, the fraction of columns made
categorical, because real tables carry categorical columns that a purely numeric SCM prior
under-produces. These are exactly the mismatches a prior designed without downstream tasks in
view would be expected to have --- and they are found, not hand-specified.

\input{knobs}

\section{Limitations}

The gains are real but modest, and they are measured on a single backbone (TabICL) at one
adapter size. We adapt rather than pre-train, so we can only \emph{repair} a prior, not
demonstrate that a better prior trained from scratch would go further; the ceiling of this
method is unknown. The agent's knowledge enters through summary statistics of the
development tasks, which means it can only correct mismatches that those statistics reveal
--- a deficiency invisible in feature counts and class balance will stay invisible. Our
search budget (\NRounds{} candidates) is small enough that random search has not converged,
which flatters the agent; with a much larger budget the two arms may meet. Finally, TabICL
does not consume column names, so the prior cannot carry semantics: the world knowledge the
LLM contributes here is knowledge about the \emph{shape} of real tabular problems, not about
what any particular column means. Injecting semantics would require a backbone that reads
headers \cite{spinaci2025contexttab}.

\section{Conclusion}

The synthetic prior of a tabular foundation model is usually treated as a fixed artefact of
its pre-training. We showed it is an editable object: exposing its generating configuration
to an LLM agent, and using LoRA to make each candidate cheap to evaluate, lets a released
checkpoint be repaired against downstream evidence for a few GPU-hours. The agent beats
random search over the same knobs under the same budget, which says the useful thing --- that
knowing what real tables look like is worth more than sampling the space uniformly.

\bibliographystyle{icml2025}
\bibliography{refs}

\end{document}
